\ifdefined\gkver\else\def\gkver{student}\fi
\documentclass[\gkver]{gaokaozhenti}
\begin{document}

\chapter{2014年海南}

\begin{enumerate}
\item
%% number: 1
%% paperName: 2014年海南高考第 1 题 3 分
%% typeId: 1
%% type: 单选题
%% chapter: 电磁感应
%% point: 楞次定律
%% method: 无
%% score: 3
%% degree: 650
%% duplicateId: 0
%% body:
如图，在一水平、固定的闭合导体圆环上方。有一条形磁铁（N 极朝上，S 极朝下）由静止开始下落，磁铁从圆环中穿过且不与圆环接触，关于圆环中感应电流的方向（从上向下看），下列说法正确的是\xzanswer{C}
\onepicture[w=3.0cm]{figs/01.png}
\fourchoices[answer=C]
{总是顺时针}
{总是逆时针}
{先顺时针后逆时针}
{先逆时针后顺时针}
%% answer: C
%% memo:
\memoanswer{
磁铁从圆环中穿过且不与圆环接触，则导体环中，先是向上的磁通量增加，磁铁过中间以后，向上的磁通量减少，根据楞次定律，产生的感应电流先顺时针后逆时针，选项 C 正确。
}

\item
%% number: 2
%% paperName: 2014年海南高考第 2 题 3 分
%% typeId: 1
%% type: 单选题
%% chapter: 交流电
%% point: 变压器
%% method: 无
%% score: 3
%% degree: 650
%% duplicateId: 0
%% body:
理想变压器上接有三个完全相同的灯泡，其中一个与该变压器的原线圈串联后接入交流电源，另外两个并联后接在副线圈两端。已知三个灯泡均正常发光。该变压器原、副线圈的匝数之比为\xzanswer{B}
\fourchoices[answer=B]
{$1:2$}
{$2:1$}
{$2:3$}
{$3:2$}
%% answer: B
%% memo:
\memoanswer{
三灯都正常工作，则电流相等，由此可知变压器的原副线圈的电流比 $\frac{I_{1}}{I_{2}}=\frac{1}{2}$，对于单匝输入单匝输出的变压器，由于功率相等，$p_{1}=p_{2}$，得 $\frac{u_{1}}{u_{2}}=\frac{I_{2}}{I_{1}}$，得 $\frac{u_{1}}{u_{2}}=\frac{2}{1}$，选项 B 正确。
}

\item
%% number: 3
%% paperName: 2014年海南高考第 3 题 3 分
%% typeId: 1
%% type: 单选题
%% chapter: 运动和力
%% point: 牛顿定律的应用
%% method: 无
%% score: 3
%% degree: 650
%% duplicateId: 0
%% body:
将一物体以某一初速度竖直上抛。物体在运动过程中受到一大小不变的空气阻力作用，它从抛出点到最高点的运动时间为 $t_{1}$，再从最高点回到抛出点的运动时间为 $t_{2}$，如果没有空气阻力作用，它从抛出点到最高点所用的时间为 $t_{0}$，则\xzanswer{B}
\fourchoices[answer=B]
{$t_{1}>t_{0}$，$t_{2}<t_{1}$}
{$t_{1}<t_{0}$，$t_{2}>t_{1}$}
{$t_{2}>t_{0}$，$t_{2}>t_{1}$}
{$t_{1}<t_{0}$，$t_{2}<t_{1}$}
%% answer: B
%% memo:
\memoanswer{
三种情况下的匀变速加速度是 $a_{1}>g>a_{2}$，其中 $a_{1}t_{1}=v_{0}=gt_{0}$，得 $t_{0}>t_{1}$；又上升与下降过程 $\frac{1}{2}a_{1}t_{1}^{2}=\frac{1}{2}a_{2}t_{2}^{2}$，得 $t_{2}>t_{1}$，选项 B 正确。
}

\item
%% number: 4
%% paperName: 2014年海南高考第 4 题 3 分
%% typeId: 1
%% type: 单选题
%% chapter: 电场
%% point: 带电粒子的直线运动
%% method: 无
%% score: 3
%% degree: 350
%% duplicateId: 0
%% body:
如图，一平行板电容器的两极板与一电压恒定的电源相连，极板水平放置，极板间距为 $d$，在下极板上叠放一厚度为 $l$ 的金属板，其上部空间有一带电粒子 P 静止在电容器中，当把金属板从电容器中快速抽出后，粒子 P 开始运动，重力加速度为 $g$。粒子运动加速度为\xzanswer{A}
\onepicture[w=5.0cm]{figs/04.png}
\fourchoices[answer=A]
{$\frac{l}{d}g$}
{$\frac{d-l}{d}g$}
{$\frac{l}{d-l}g$}
{$\frac{d}{d-l}g$}
%% answer: A
%% memo:
\memoanswer{
平行板电容器的两极板与一电压恒定的电源相连，有金属板时，板间场强可以表达为 $E_{1}=\frac{U}{d-l}$，且有 $E_{1}q=mg$；当抽去金属板，则板间距离增大，板间场强可以表达为 $E_{2}=\frac{U}{d}$，有 $mg-E_{2}q=ma$，联立上述可解得 $\frac{d}{-l}=\frac{g}{-a}$，知选项 A 正确。
}

\item
%% number: 5
%% paperName: 2014年海南高考第 5 题 3 分
%% typeId: 1
%% type: 单选题
%% chapter: 力物体的平衡
%% point: 三力平衡
%% method: 无
%% score: 3
%% degree: 650
%% duplicateId: 0
%% body:
如图，一不可伸长的光滑轻绳，其左端固定于 $O$ 点，右端跨过位于 $O'$ 点的固定光滑轴悬挂一质量为 $M$ 的物体；$OO'$ 段水平，长度为 $L$；绳子上套一可沿绳滑动的轻环。现在轻环上悬挂一钩码，平衡后，物体上升 $L$。则钩码的质量为\xzanswer{D}
\onepicture[w=3.6cm]{figs/05.png}
\fourchoices[answer=D]
{$\frac{\sqrt{2}}{2}M$}
{$\frac{\sqrt{3}}{2}M$}
{$\sqrt{2}M$}
{$\sqrt{3}M$}
%% answer: D
%% memo:
\memoanswer{
平衡后，物体上升 $L$，说明环下移后，将绳子拉过来的长度为 $L$，取环重新平衡的位置为 A 点，则 $OA=O'A=L$。

\onepicture[w=3.4cm]{figs/05_an.png}

则如图易知 $mg=\sqrt{3}Mg$，选项 D 正确。
}

\item
%% number: 6
%% paperName: 2014年海南高考第 6 题 3 分
%% typeId: 1
%% type: 单选题
%% chapter: 曲线运动
%% point: 万有引力定律
%% method: 无
%% score: 3
%% degree: 450
%% duplicateId: 0
%% body:
设地球自转周期为 $T$，质量为 $M$，引力常量为 $G$。假设地球可视为质量均匀分布的球体，半径为 $R$。同一物体在南极和赤道水平面上静止时所受到的支持力之比为\xzanswer{A}
\fourchoices[answer=A]
{$\frac{GMT^{2}}{GMT^{2}-4\pi^{2}R^{3}}$}
{$\frac{GMT^{2}}{GMT^{2}+4\pi^{2}R^{3}}$}
{$\frac{GMT^{2}-4\pi^{2}R^{3}}{GMT^{2}}$}
{$\frac{GMT^{2}+4\pi^{2}R^{3}}{GMT^{2}}$}
%% answer: A
%% memo:
\memoanswer{
物体在南极地面所受的支持力等于万有引力，$F=G\frac{Mm}{R^{2}}$；在赤道处，$F_{\text{万}}-F'=F_{\text{向}}$，得 $F'=F_{\text{万}}-F_{\text{向}}$；又 $F_{\text{向}}=m\frac{4\pi^{2}}{T^{2}}R$，则 $F'=G\frac{Mm}{R^{2}}-m\frac{4\pi^{2}}{T^{2}}R$，联立得 $\frac{F}{F'}=\frac{GMT^{2}}{GMT^{2}-4\pi^{2}R^{3}}$，A 正确。
}

\item
%% number: 7
%% paperName: 2014年海南高考第 7 题 5 分
%% typeId: 2
%% type: 多选题
%% chapter: 其他
%% point: 物理学史
%% method: 无
%% score: 5
%% degree: 650
%% duplicateId: 0
%% body:
下列说法中，符合物理学史实的是\xzanswer{ABD}
\fourchoices[answer=ABD]
{亚里士多德认为，必须有力作用在物体上，物体才能运动；没有力的作用，物体就静止}
{牛顿认为，力是物体运动状态改变的原因，而不是物体运动的原因}
{麦克斯韦发现了电流的磁效应，即电流可以在其周围产生磁场}
{奥斯特发现导线通电时，导线附近的小磁针发生偏转}
%% answer: ABD
%% memo:
\memoanswer{
在亚里士多德的时代，人们根据生产生活的实际经验感受到，物体的运动都是需要有力的作用的，A 正确；牛顿根据伽利略的研究成果提出惯性的概念，并继承和发展了伽利略的观点，认为物体运动的维持不需要力，运动的改变才需要力，B 正确；奥斯特发现通电导线下方的小磁针会偏转，从而发现了电流磁效应，故 C 错 D 正确。
}

\item
%% number: 8
%% paperName: 2014年海南高考第 8 题 5 分
%% typeId: 2
%% type: 多选题
%% chapter: 磁场
%% point: 磁场的叠加
%% method: 无
%% score: 5
%% degree: 650
%% duplicateId: 0
%% body:
如图，两根平行长直导线相距 $2L$，通有大小相等、方向相同的恒定电流；a、b、c 是导线所在平面内的三点，左侧导线与它们的距离分别为 $\frac{l}{2}$、$l$ 和 $3l$。关于这三点处的磁感应强度，下列判断正确的是\xzanswer{AD}
\onepicture[w=3.4cm]{figs/08.png}
\fourchoices[answer=AD]
{a 处的磁感应强度大小比 c 处的大}
{b、c 两处的磁感应强度大小相等}
{a、c 两处的磁感应强度方向相同}
{b 处的磁感应强度为零}
%% answer: AD
%% memo:
\memoanswer{
根据通电直导线的磁场，利用右手螺旋定则，可知 b 处场强为零，两导线分别在 a 处的产生的场强都大于在 c 处产生的场强，a、c 两处的场强叠加都是同向叠加，选项 AD 正确。
}

\item
%% number: 9
%% paperName: 2014年海南高考第 9 题 5 分
%% typeId: 2
%% type: 多选题
%% chapter: 电场
%% point: 带电粒子的直线运动
%% method: 无
%% score: 5
%% degree: 450
%% duplicateId: 0
%% body:
如图（a），直线 $MN$ 表示某电场中一条电场线，a、b 是线上的两点，将一带负电荷的粒子从 a 点处由静止释放，粒子从 a 运动到 b 过程中的 $v-t$ 图线如图（b）所示，设 a、b 两点的电势分别为 $\varphi_{a}$、$\varphi_{b}$，场强大小分别为 $E_{a}$、$E_{b}$，粒子在 a、b 两点的电势能分别为 $W_{a}$、$W_{b}$，不计重力，则有\xzanswer{BD}
\onepicture[w=6.5cm]{figs/09.png}
\fourchoices[answer=BD]
{$\varphi_{a}>\varphi_{b}$}
{$E_{a}>E_{b}$}
{$E_{a}<E_{b}$}
{$W_{a}>W_{b}$}
%% answer: BD
%% memo:
\memoanswer{
由 $v-t$ 图像的斜率渐小可知由 a 到 b 的过程中，粒子的加速度渐小，所以场强渐小，$E_{a}>E_{b}$；根据动能定理，速度增大，可知势能减小，$W_{a}>W_{b}$，可得选项 BD 正确。
}

\item
%% number: 10
%% paperName: 2014年海南高考第 10 题 5 分
%% typeId: 2
%% type: 多选题
%% chapter: 功和能
%% point: 功能原理
%% method: 无
%% score: 5
%% degree: 450
%% duplicateId: 0
%% body:
如图，质量相同的两物体 a、b，用不可伸长的轻绳跨接在同一光滑的轻质定滑轮两侧，a 在水平桌面的上方，b 在水平粗糙桌面上。初始时用力压住 b 使 a、b 静止，撤去此压力后，a 开始运动，在 a 下降的过程中，b 始终未离开桌面。在此过程中\xzanswer{AD}
\onepicture[w=4.5cm]{figs/10.png}
\fourchoices[answer=AD]
{a 的动能小于 b 的动能}
{两物体机械能的变化量相等}
{a 的重力势能的减小量等于两物体总动能的增加量}
{绳的拉力对 a 所做的功与对 b 所做的功的代数和为零}
%% answer: AD
%% memo:
\memoanswer{
由于 $v_{a}=v_{b}\cos\theta$，$\theta$ 为 b 的拉绳与水平面的夹角，质量相同，动能 $E_{k}=\frac{1}{2}mv^{2}$，可知选项 A 正确；a 物体下降时，a 的机械能的减少量等于 b 物体的动能增加量和 b 克服摩擦力做功之和，选项 BC 错误；绳的拉力对 a 所做的功等于 a 的机械能的减少量，绳的拉力对 b 所做的功等于 b 的动能增加和克服摩擦力做功，选项 D 正确。
}

\item
%% number: 11
%% paperName: 2014年海南高考第 11 题 5 分
%% typeId: 4
%% type: 实验题
%% chapter: 电路
%% point: 电阻定律
%% method: 无
%% score: 5
%% degree: 650
%% duplicateId: 0
%% body:
现有一合金制成的圆柱体。为测量该合金的电阻率，现用伏安法测圆柱体两端之间的电阻，用螺旋测微器测量该圆柱体的直径，用游标卡尺测量该圆柱体的长度。螺旋测微器和游标卡尺的示数如图（a）和图（b）所示。
\begin{enumerate}
	\item
	由上图读得圆柱体的直径为\tkanswer[2cm]{1.845}$\Umm$，长度为\tkanswer[2cm]{4.240}$\Ucm$。
	\item
	若流经圆柱体的电流为 $I$，圆柱体两端之间的电压为 $U$，圆柱体的直径和长度分别为 $D$、$L$，测得 $D$、$L$、$I$、$U$ 表示的电阻率的关系式为 $\rho=$\tkanswer[2.5cm]{$\frac{\pi D^{2}U}{4IL}$}。
\end{enumerate}
\onepicture[align=right, w=8.0cm]{figs/11.png}
%% answer:
\jdanswer{
\begin{enumerate}
	\item
	$1.845\Umm$，$4.240\Ucm$。
	\item
	$\frac{\pi D^{2}U}{4IL}$。
\end{enumerate}
}
%% memo:
\memoanswer{
（1）螺旋测微器读数：$1.5\Umm+34.5\times0.01\Umm=1.845\Umm$，游标卡尺读数：$42\Umm+8\times0.05\Umm=42.40\Umm$。

（2）电阻 $R=\rho\frac{l}{s}$，得 $\rho=\frac{Rs}{l}$，可得 $\rho=\frac{\pi D^{2}U}{4IL}$。
}

\item
%% number: 12
%% paperName: 2014年海南高考第 12 题 10 分
%% typeId: 4
%% type: 实验题
%% chapter: 电路
%% point: 伏安法测电阻
%% method: 无
%% score: 10
%% degree: 450
%% duplicateId: 0
%% body:
用伏安法测量一电池的内阻。已知该待测电池的电动势 $E$ 约为 $9\UV$，内阻约数十欧，允许输出的最大电流为 $50\UmA$，可选用的实验器材有：

电压表 V$_{1}$（量程 $5\UV$）；电压表 V$_{2}$（量程 $10\UV$）；

电流表 A$_{1}$（量程 $50\UmA$）；电压表 A$_{2}$（量程 $100\UmA$）；

滑动变阻器 $R$（最大电阻 $300\UO$）；

定值电阻 $R_{1}$（阻值为 $200\UO$，额定功率为 $1/8\UW$）；定值电阻 $R_{2}$（阻值为 $220\UO$，额定功率为 $1\UW$）；

开关 S；导线若干。

测量数据如坐标纸上 $U-I$ 图线所示。
\begin{enumerate}
	\item
	在答题卡相应的虚线方框内画出合理的电路原理图，并标明所选器材的符号；
	\item
	在设计的电路中，选择定值电阻的根据是\_\_\_\_\_\_\_\_\_\_；
	\item
	由 $U-I$ 图线求得待测电池的内阻为\tkanswer[2cm]{51.0}$\UO$；
	\item
	在你设计的电路中，产生系统误差的主要原因是\_\_\_\_\_\_\_\_\_。
\end{enumerate}
\onepicture[align=right, w=7.0cm]{figs/12.png}
%% answer:
\jdanswer{
\begin{enumerate}
	\item
	如图（a）所示（给出图（b）也给分）。
	\item
	定值电阻在电路中消耗的功率会超过 $1/8\UW$，$R_{2}$ 的功率满足实验要求。
	\item
	51.0。
	\item
	忽略了电压表的分流（此答案对应于图（a）），或忽略了电流表的分压（此答案对应于图（b））。
\end{enumerate}
}
%% memo:
\memoanswer{
（1）电路原理图如图（a）所示。（5 分，给出图（b）也给分。原理正确 2 分，仪器选择正确 3 分）

\onepicture[w=9.0cm]{figs/12_an.png}

（2）定值电阻在电路中消耗的功率会超过 $1/8\UW$，$R_{2}$ 的功率满足实验要求（1 分）。

（3）$51.0\UO$（2 分。在 $49.0\sim53.0\UO$ 范围内的均给分）。

（4）忽略了电压表的分流（此答案对应于图（a））或忽略了电流表的分压（此答案对应于图（b））（2 分，其他合理答案也给分）。
}

\item
%% number: 13
%% paperName: 2014年海南高考第 13 题 9 分
%% typeId: 6
%% type: 计算题
%% chapter: 直线运动
%% point: 多段匀变速直线运动
%% method: 无
%% score: 9
%% degree: 650
%% duplicateId: 0
%% body:
短跑运动员完成 $100\Um$ 赛跑的过程可简化为匀加速直线运动和匀速直线运动两个阶段。一次比赛中，某运动员用 $11.00\Us$ 跑完全程。已知运动员在加速阶段的第 $2\Us$ 内通过的距离为 $7.5\Um$，求该运动员的加速度及在加速阶段通过的距离。
%% answer:
\jdanswer{
$a=5\Umsq$，$s'=10\Um$。
}
%% memo:
\memoanswer{
根据题意，在第 1 s 和第 2 s 内运动员都做匀加速直线运动，设运动员在匀加速阶段的加速度为 $a$，在第 1 s 和第 2 s 内通过的位移分别为 $s_{1}$ 和 $s_{2}$，由运动学规律得 $s_{1}=\frac{1}{2}at_{0}^{2}$，$s_{1}+s_{2}=\frac{1}{2}a(2t_{0})^{2}$，$t_{0}=1\Us$，求得 $a=5\Umsq$。

设运动员做匀加速运动的时间为 $t_{1}$，匀速运动的时间为 $t_{2}$，匀速运动的速度为 $v_{1}$，跑完全程的时间为 $t$，全程的距离为 $s$，依题意及运动学规律，得 $t=t_{1}+t_{2}$，$v=at_{1}$，$s=\frac{1}{2}at_{1}^{2}+vt_{2}$。

设加速阶段通过的距离为 $s'$，则 $s'=\frac{1}{2}at_{1}^{2}$，求得 $s'=10\Um$。
}

\item
%% number: 14
%% paperName: 2014年海南高考第 14 题 14 分
%% typeId: 6
%% type: 计算题
%% chapter: 磁场
%% point: 带电粒子在磁场中的运动
%% method: 无
%% score: 14
%% degree: 450
%% duplicateId: 0
%% body:
如图，在 $x$ 轴上方存在匀强磁场，磁感应强度大小为 $B$，方向垂直于纸面向外；在 $x$ 轴下方存在匀强电场，电场方向与 $xOy$ 平面平行，且与 $x$ 轴成 $45^\circ$ 夹角。一质量为 $m$、电荷量为 $q$（$q>0$）的粒子以速度 $v_{0}$ 从 $y$ 轴上 P 点沿 $y$ 轴正方向射出，一段时间后进入电场，进入电场时的速度方向与电场方向相反；又经过一段时间 $T_{0}$，磁场方向变为垂直纸面向里，大小不变，不计重力。
\begin{enumerate}
	\item
	求粒子从 P 点出发至第一次到达 $x$ 轴时所需的时间；
	\item
	若要使粒子能够回到 P 点，求电场强度的最大值。
\end{enumerate}
\onepicture[align=right, w=5.0cm]{figs/14.png}
%% answer:
\jdanswer{
\begin{enumerate}
	\item
	$\frac{5\pi m}{4qB}$
	\item
	$\frac{2mv_{0}}{qT_{0}}$
\end{enumerate}
}
%% memo:
\memoanswer{
（1）带电粒子在磁场中做圆周运动，设运动半径为 $R$，运动周期为 $T$，根据洛伦兹力公式及圆周运动规律，有 $qv_{0}B=m\frac{v_{0}^{2}}{R}$，$T=\frac{2\pi R}{v_{0}}$。依题意，粒子第一次到达 $x$ 轴时，运动转过的角度为 $\frac{5}{4}\pi$，所需时间 $t_{1}$ 为 $t_{1}=\frac{5}{8}T$，求得 $t_{1}=\frac{5\pi m}{4qB}$。

（2）粒子进入电场后，先做匀减速运动，直到速度减小为 0，然后沿原路返回做匀加速运动，到达 $x$ 轴时速度大小仍为 $v_{0}$，设粒子在电场中运动的总时间为 $t_{2}$，加速度大小为 $a$，电场强度大小为 $E$，有 $qE=ma$，$v_{0}=\frac{1}{2}at_{2}$，得 $t_{2}=\frac{2mv_{0}}{qE}$。根据题意，要使粒子能够回到 P 点，必须满足 $t_{2}\geq T_{0}$，得电场强度最大值 $E=\frac{2mv_{0}}{qT_{0}}$。
}

\item
%% number: 15
%% paperName: 2014年海南高考第 15 题 8 分
%% typeId: 6
%% type: 计算题
%% chapter: 气体性质
%% point: 气体连接体
%% method: 无
%% score: 8
%% degree: 650
%% duplicateId: 0
%% body:
模块 3-3 试题（12 分）
\begin{enumerate}
	\item
	（4 分）下列说法正确的是\tkanswer[1.6cm]{CE}。（填入正确答案标号，选对 1 个得 2 分，选对 2 个得 4 分；有选错的得 0 分）
	\fivechoices[answer=CE]
	{液体表面张力的方向与液面垂直并指向液体内部}
	{单晶体有固定的熔点，多晶体没有固定的熔点}
	{单晶体中原子（或分子、离子）的排列具有空间周期性}
	{通常金属在各个方向的物理性质都相同，所以金属是非晶体}
	{液晶具有液体的流动性，同时具有晶体的各向异性特征}
	\item
	（8 分）一竖直放置、缸壁光滑且导热的柱形气缸内盛有一定量的氮气，被活塞分隔成Ⅰ、Ⅱ两部分；达到平衡时，这两部分气体的体积相等，上部气体的压强为 $p_{10}$，如图（a）所示，若将气缸缓慢倒置，再次达到平衡时，上下两部分气体的体积之比为 $3:1$，如图（b）所示。设外界温度不变，已知活塞面积为 $S$，重力加速度大小为 $g$，求活塞的质量。
	\onepicture[w=3.6cm, align=right]{figs/15.png}
\end{enumerate}
%% answer:
\jdanswer{
\begin{enumerate}
	\item
	CE。
	\item
	$\frac{4p_{10}S}{5g}$。
\end{enumerate}
}
%% memo:
\memoanswer{
（1）液体表面张力的方向始终与液面相切，A 错误；单晶体和多晶体都有固定的熔沸点，非晶体熔点不固定，B 错误；单晶体中原子（或分子、离子）的排列是规则的，具有空间周期性，表现为各向异性，C 正确；金属材料虽然显示各向同性，但并不意味着就是非晶体，可能是多晶体，D 错误；液晶的名称由来就是由于它具有流动性和各向异性，故 E 正确。

（2）设活塞的质量为 $m$，气缸倒置前下部气体的压强为 $p_{20}$，倒置后上下气体的压强分别为 $p_{2}$、$p_{1}$，由力的平衡条件有 $p_{20}=p_{10}+\frac{mg}{S}$，$p_{1}=p_{2}+\frac{mg}{S}$。倒置过程中，两部分气体均经历等温过程，设气体的总体积为 $V_{0}$，由玻意耳定律得 $p_{10}\frac{V_{0}}{2}=p_{1}\frac{V_{0}}{4}$，$p_{20}\frac{V_{0}}{2}=p_{2}\frac{V_{0}}{4}$，解得 $m=\frac{4p_{10}S}{5g}$。
}

\item
%% number: 16
%% paperName: 2014年海南高考第 16 题 8 分
%% typeId: 6
%% type: 计算题
%% chapter: 光学
%% point: 全反射
%% method: 无
%% score: 8
%% degree: 650
%% duplicateId: 0
%% body:
模块 3-4 试题（12 分）
\begin{enumerate}
	\item
	（4 分）一列简谐横波沿 $x$ 轴传播，a、b 为 $x$ 轴上的两质点，平衡位置分别为 $x=0$，$x=x_{b}$（$x_{b}>0$）。a 点的振动规律如图所示，已知波速为 $v=10\Ums$，在 $t=0.1\Us$ 时，b 点的位移为 $0.05\Um$，则下列判断可能正确的是\tkanswer[1.6cm]{BC}。（填入正确答案标号，选对 1 个得 2 分，选对 2 个得 4 分；有选错的得 0 分）
	\fourchoices[answer=BC]
	{波沿 $x$ 轴正向传播，$x_{b}=0.5\Um$}
	{波沿 $x$ 轴正向传播，$x_{b}=1.5\Um$}
	{波沿 $x$ 轴负向传播，$x_{b}=2.5\Um$}
	{波沿 $x$ 轴负向传播，$x_{b}=3.5\Um$}
	\onepicture[w=6.0cm, align=right]{figs/16a.png}
	\item
	（8 分）如图，矩形 $ABCD$ 为一水平放置的玻璃砖的截面，在截面所在平面有一细束激光照射玻璃砖，入射点距底面的高度为 $h$，反射光线和折射光线与底面所在平面的交点到 AB 的距离分别 $l_{1}$ 和 $l_{2}$，在截面所在平面内，改变激光束在 AB 面上入射点的高度与入射角的大小，当折射光线与底面的交点到 AB 的距离为 $l_{3}$ 时，光线恰好不能从底面射出，求此时入射点距离底面的高度 $H$。
	\onepicture[w=4.5cm, align=right]{figs/16b.png}
\end{enumerate}
%% answer:
\jdanswer{
\begin{enumerate}
	\item
	BC。
	\item
	$\sqrt{\frac{l_{2}^{2}-l_{1}^{2}}{l_{1}^{2}+h^{2}}}\,l_{3}$。
\end{enumerate}
}
%% memo:
\memoanswer{
（1）由题知，该波波长为 $\lambda=vT=2\Um$，若波沿 $x$ 轴正向传播，且 $x_{b}=1.5\Um$，则 a、b 距离为四分之三个波长，当 a 在平衡位置向上振动时，b 刚好在正向最大位移处，即位移为 $0.05\Um$，B 正确；若波沿 $x$ 轴负向传播，且 $x_{b}=2.5\Um$（一又四分之一波长处），当 a 在平衡位置向上振动时 b 刚好也处在正向最大位移处，即位移为 $0.05\Um$，C 正确。

（2）设玻璃砖的折射率为 $n$，入射角和反射角为 $\theta_{1}$，折射角为 $\theta_{2}$，由光的折射定律 $n=\frac{\sin\theta_{1}}{\sin\theta_{2}}$。根据几何关系有 $\sin\theta_{1}=\frac{h}{\sqrt{l_{1}^{2}+h^{2}}}$，$\sin\theta_{2}=\frac{h}{\sqrt{l_{2}^{2}+h^{2}}}$，因此求得 $n=\sqrt{\frac{l_{2}^{2}+h^{2}}{l_{1}^{2}+h^{2}}}$。根据题意，折射光线在某一点刚好无法从底面射出，此时发生全反射，设在底面发生全反射时的入射角为 $\theta_{3}$，有 $\sin\theta_{3}=\frac{1}{n}$，由几何关系得 $\sin\theta_{3}=\frac{l_{3}}{\sqrt{l_{3}^{2}+H^{2}}}$，解得 $H=\sqrt{\frac{l_{2}^{2}-l_{1}^{2}}{l_{1}^{2}+h^{2}}}\,l_{3}$。
}

\item
%% number: 17
%% paperName: 2014年海南高考第 17 题 8 分
%% typeId: 6
%% type: 计算题
%% chapter: 运动和力
%% point: 动量守恒定律
%% method: 无
%% score: 8
%% degree: 450
%% duplicateId: 0
%% body:
模块 3-5 试题（12 分）
\begin{enumerate}
	\item
	（4 分）在光电效应实验中，用同一种单色光，先后照射锌和银的表面，都能发生光电效应。对于这两个过程，下列四个物理过程中，一定不同的是\tkanswer[1.6cm]{ACD}。（填入正确答案标号，选对 1 个得 2 分，选对 2 个得 4 分；有选错的得 0 分）
	\fourchoices[answer=ACD]
	{遏止电压}
	{饱和光电流}
	{光电子的最大初动能}
	{逸出功}
	\item
	（8 分）一静止原子核发生 $\alpha$ 衰变，生成一 $\alpha$ 粒子及一新核，$\alpha$ 粒子垂直进入磁感应强度大小为 $B$ 的匀强磁场，其运动轨迹是半径为 $R$ 的圆。已知 $\alpha$ 粒子的质量为 $m$，电荷量为 $q$；新核的质量为 $M$；光在真空中的速度大小为 $c$。求衰变前原子核的质量。
\end{enumerate}
%% answer:
\jdanswer{
\begin{enumerate}
	\item
	ACD。
	\item
	$(M+m)\left[1+\frac{(qBR)^{2}}{2Mmc^{2}}\right]$。
\end{enumerate}
}
%% memo:
\memoanswer{
（1）不同的金属具有不同的逸出功，遏止电压为 $U=\frac{h\nu-W_{0}}{e}$，光电子的最大初动能为 $E_{k}=h\nu-W_{0}$，饱和光电流由单位时间内的入射光子数决定，综上可知 ACD 正确。

（2）设衰变产生的 $\alpha$ 粒子的速度大小为 $v$，由洛伦兹力公式和牛顿第二定律得 $qvB=m\frac{v^{2}}{R}$。设衰变后新核的速度大小为 $V$，衰变前后动量守恒，有 $0=MV-mv$。设衰变前原子核质量为 $M_{0}$，衰变前后能量守恒，有 $M_{0}c^{2}=Mc^{2}+\frac{1}{2}MV^{2}+mc^{2}+\frac{1}{2}mv^{2}$，解得 $M_{0}=(M+m)\left[1+\frac{(qBR)^{2}}{2Mmc^{2}}\right]$。
}

\end{enumerate}

\end{document}
